\section{Profundización teórica: una perspectiva unificada desde DDIM hasta DPM-Solver}

\subsection{Re-derivación de la fórmula de muestreo de DDIM}

El docente llevó a los estudiantes a expandir término por término los resultados de DPM-Solver-1 para verificar que son equivalentes a DDIM:

Partiendo de la fórmula de DPM-Solver-1:
$$\mathbf{x}_t = \frac{\alpha_t}{\alpha_s} \mathbf{x}_s - \sigma_t (e^{h} - 1) \boldsymbol{\epsilon}_\theta(\mathbf{x}_s, s)$$

donde $h = \lambda_t - \lambda_s = \log\frac{\alpha_t}{\sigma_t} - \log\frac{\alpha_s}{\sigma_s} = \log\frac{\alpha_t \sigma_s}{\sigma_t \alpha_s}$.

Por tanto, $e^h = \frac{\alpha_t \sigma_s}{\sigma_t \alpha_s}$ y $e^h - 1 = \frac{\alpha_t \sigma_s - \sigma_t \alpha_s}{\sigma_t \alpha_s}$.

Sustituyendo:
$$\mathbf{x}_t = \frac{\alpha_t}{\alpha_s} \mathbf{x}_s - \sigma_t \cdot \frac{\alpha_t \sigma_s - \sigma_t \alpha_s}{\sigma_t \alpha_s} \cdot \boldsymbol{\epsilon}_\theta(\mathbf{x}_s, s)$$
$$= \frac{\alpha_t}{\alpha_s} \mathbf{x}_s - \frac{\alpha_t \sigma_s - \sigma_t \alpha_s}{\alpha_s} \cdot \boldsymbol{\epsilon}_\theta(\mathbf{x}_s, s)$$

\begin{importantbox}{Perspectiva unificada}
Esta es precisamente la fórmula de muestreo de DDIM bajo parametrización de tiempo continuo. A través de esta derivación, observamos las conexiones profundas entre diferentes formas de entender los modelos de difusión:
\begin{enumerate}
    \item DDPM $\to$ DDIM (Song et al., 2020): de muestreo estocástico a muestreo determinista
    \item SDE $\to$ PF-ODE (Song et al., 2021): de proceso estocástico a flujo determinista
    \item PF-ODE $\to$ integral exponencial (Lu et al., 2022): de método numérico general a solucionador estructurado
\end{enumerate}
DDIM se sitúa exactamente en la posición de aproximación de primer orden de la última línea: no fue "inventado", sino que es un producto natural de la estructura ODE de los modelos de difusión.
\end{importantbox}

\subsection{Desarrollo de Taylor en el espacio $\lambda$}

El marco unificado de los métodos de varios órdenes de DPM-Solver puede entenderse mediante el desarrollo de Taylor de $\boldsymbol{\epsilon}_\theta$ en el espacio $\lambda$:
$$\boldsymbol{\epsilon}_\theta(\lambda) = \boldsymbol{\epsilon}_\theta(\lambda_s) + (\lambda - \lambda_s) \boldsymbol{\epsilon}_\theta'(\lambda_s) + \frac{(\lambda - \lambda_s)^2}{2} \boldsymbol{\epsilon}_\theta''(\lambda_s) + \cdots$$

\begin{itemize}
    \item \textbf{Corte de orden cero} ($\boldsymbol{\epsilon}_\theta \approx$ constante) $\to$ DPM-Solver-1 = DDIM
    \item \textbf{Corte de primer orden} ($\boldsymbol{\epsilon}_\theta \approx$ lineal) $\to$ DPM-Solver-2
    \item \textbf{Corte de segundo orden} ($\boldsymbol{\epsilon}_\theta \approx$ cuadrático) $\to$ DPM-Solver-3
\end{itemize}

Las derivadas de alto orden $\boldsymbol{\epsilon}_\theta'(\lambda_s)$ no pueden obtenerse analíticamente (ya que $\boldsymbol{\epsilon}_\theta$ es una red neuronal), pero pueden aproximarse mediante diferencias finitas: esta es la razón por la cual es necesario evaluar adicionalmente la red en puntos intermedios.

\subsection{Resumen de la sección}

Desde la perspectiva del desarrollo de Taylor, los métodos de varios órdenes de DPM-Solver forman una familia natural de precisión creciente. Se revela la identidad de DDIM como caso particular de primer orden, mientras que los métodos de alto orden realizan un ajuste polinomial más refinado de la predicción de ruido en el espacio $\lambda$.

